\originalchapter{研究生作业十：推迟场与经典辐射}
题面据8.311作业10全部5题, 解答为AI独立推导并与17页官方解答逐项对照. 所谓经典氢轨道仅为理论模型, 不把经典塌缩解释成真实量子氢原子的寿命.
\originalsection{匀速电荷的势与场}
\begin{exercise}
原题1. 点电荷$q$匀速$\vct v$, (a) 将静止Coulomb四势变换到运动系; (b) 微分求场并与匀速LW场核对.
\end{exercise}
\begin{solution}
静止势$V'=q/(4\pi\eps r')$, $\vct A'=0$. 反boost给$V=\gamma V'$、$\vct A=\vct v V/c^2$, 同时$r'^2=\gamma^2R_{0\parallel}^2+R_{0\perp}^2$, $\vct R_0=\rr-\vct vt$（初始位置可平移）. 因而$V=q/[4\pi\eps\sqrt{R_0^2-|\vct\beta\times\vct R_0|^2}]$.

记$D^2=R_0^2-(\beta^2R_0^2-(\vct\beta\cdot\vct R_0)^2)$. $\nabla D^2=2[(1-\beta^2)\vct R_0+(\vct\beta\cdot\vct R_0)\vct\beta]$, $\partial_tD^2=-2\vct v\cdot\vct R_0$. 在$-\nabla V-\vct v\partial_tV/c^2$中平行$\vct\beta$项抵消, 得$\ee=q(1-\beta^2)\vct R_0/(4\pi\eps D^3)$, $\bb=\vct v\times\ee/c^2$. 此场指向瞬时位置而大小具横向增强, 没有$1/r$辐射项; 与第7章$\dot\beta=0$的迟滞表示等价.
\end{solution}
\originalsection{LW势的根与全场求导}
\begin{exercise}
原题2. 从点源推迟积分推导Liénard--Wiechert势, 再实际微分得任意亚光速运动电荷全场.
\end{exercise}
\begin{solution}
本题主计算链已在第7章“点电荷的推迟势”及该章新增的完整求导段展开, 此处定位各步骤以供核对: 源积分的$h(t')=t-t'-|\rr-\rr_q(t')|/c$有$h'=-\kappa$, $\kappa=1-\nn\cdot\vct\beta>0$, 故三维delta及时间根给$V=q/(4\pi\eps\kappa R)$、$\vct A=\vct vV/c^2$. 对观察变量$\partial_tt_r=1/\kappa$, $\nabla t_r=-\nn/(c\kappa)$, 设$F=\kappa R$, 第7章已逐项计算$\nabla F,\partial_tF$并合并两势. 最终
\[
\ee=\frac q{4\pi\eps}\left[\frac{(1-\beta^2)(\nn-\vct\beta)}{\kappa^3R^2}+\frac{\nn\times[(\nn-\vct\beta)\times\dot{\vct\beta}]}{c\kappa^3R}\right]_{t_r},\qquad\bb=\nn\times\ee/c.
\]
各量同取迟滞时刻. 第7章的势梯度与矢势时间导数逐项相加是本题的完整中文推导, 此处不把仅列结果当另一次证明. 同一卷内复用可减少重复, 并保留$\kappa^{-3}$由根导数与势导数共同产生的解释.
\end{solution}
\originalsection{非相对论回旋能量衰减}
\begin{exercise}
原题3. 电子在恒磁场做垂直圆运动, 非相对论且慢辐射损失. 求$P$, 导能量衰减时间并对$10^4$gauss估算.
\end{exercise}
\begin{solution}
$\Omega=eB/m$, $a=\Omega v$, $P=e^2\Omega^2v^2/(6\pi\eps c^3)$. 动能$K=mv^2/2$使$\dot K=-\Gamma K$, $\Gamma=e^4B^2/(3\pi\eps m^3c^3)=2\tau_e\Omega^2$, $\tau_e=e^2/(6\pi\eps mc^3)$. 故$K(t)=K_0e^{-t/\tau_{
\rm loss}}$, $\tau_{\rm loss}=1/\Gamma$. $10^4$gauss$=1$T, 电子给$\tau_{\rm loss}\simeq2.58$s, 与官方约2.59s一致. 磁场不做功, 半径缓慢缩小; 采用$\Gamma\ll\Omega$而非恒半径无限辐射.
\end{solution}
\originalsection{经典氢的有限塌缩模型}
\begin{exercise}
原题4. 带电$q$粒子绕静止$-q$近圆运动并辐射, $\kappa=q^2/(4\pi\eps)$. 求能量从$E_2$变为$E_1<E_2<0$的时长, 证模型有限塌缩, 以$r_0=10^{-8}$cm估计电子时长.
\end{exercise}
\begin{solution}
圆轨道$v^2=\kappa/(mr)$, $E=-\kappa/(2r)$, $a=\kappa/(mr^2)$, Larmor给$P=2\kappa^3/(3m^2c^3r^4)=AE^4$, $A=32/(3m^2c^3\kappa)$. 因$\dot E=-AE^4$, 积分得到$\Delta t=(E_1^{-3}-E_2^{-3})/(3A)>0$. 正确符号可由初值$E_2$较不负检查.

形式上$E_1\to-\infty$仍给有限$t_{\rm col}=-1/(3AE_2^3)=m^2c^3r_0^3/(4\kappa^2)$. $r_0=10^{-10}$m、电子电荷和质量给$1.05\times10^{-10}$s. 近中心$v$增大, 非相对论和慢变近圆假定失效, 点源自身问题亦出现; 因而有限塌缩是该经典近似的矛盾征兆, 不是在失效区域仍精确的预测.
\end{solution}
\originalsection{直线驱动与回旋角分布}
\begin{exercise}
原题5. (a) 弱简谐电场$E_0\cos\omega t$驱动自由电子, 求平均角辐射; (b) 电子非相对论磁回旋, 求周期平均角分布及总功率.
\end{exercise}
\begin{solution}
$\ddot{\pp}=q\vct a$把粒子加速度与偶极辐射工具连接. (a) 沿电场加速度$a=-eE_0\cos\omega t/m$, 观测方向与该轴夹角$\theta$, 平均$\langle a^2\rangle=e^2E_0^2/(2m^2)$. 所以$\dd\langle P\rangle/\dd\Omega=e^4E_0^2\sin^2\theta/(32\pi^2\eps m^2c^3)$. 对初始位置零可加位置常量, 不影响加速度. 弱波忽略波自身磁力的条件$v\ll c$明确保留.

(b) 加速度在$xy$旋转、模$a=v\Omega$不变, $\langle(\nn\cdot\vct a)^2\rangle=a^2\sin^2\theta/2$, $\theta$从磁轴量. 因而$\dd\langle P\rangle/\dd\Omega=e^2a^2(1+\cos^2\theta)/(32\pi^2\eps c^3)$. 积分$\int(1+\cos^2\theta)\dd\Omega=16\pi/3$给$P=e^2a^2/(6\pi\eps c^3)$, 与题3一致. 直线振子的轴上零辐射与圆周运动的轴上非零不同.
\end{solution}
